\chapter{Discussion and conclusions}
\label{ch:discussion}
\section{Which parts of the system work?}
The clearest outcome is a concrete data product. A request specifies an initial local coordinate, a participant class and the network. A trained diffusion model samples timed movement and an explicit decoder returns positions. Across the complete final pools the system produced 1,272 trajectories, each containing 512 regularly spaced observations. These data can be read by a downstream pipeline using the same five fields as the canonical input.

The geographic result has two parts. Once a request is admitted, planned motion is constrained by the path representation and passes the current road audit in the experiment. Before generation, the strict eligibility rule removes a substantial fraction of the held-out IDs. Presenting only the first result would obscure the largest practical limitation. A system serving arbitrary start-class pairs or arbitrary traffic states would have to account for trajectories excluded because of duration, lateral movement, reconstruction or participant type.

The matched empirical baseline identifies a narrower contribution of diffusion. Relative to transferring observed motion to identical routes, diffusion improved the class-averaged travel-distance distribution on the locked test and avoided exact time-aligned reuse of whole training profiles. Speed and acceleration improvements were uncertain, and the empirical samples exhibited greater spatial separation. These results support modelling the motion law on top of the road structure while limiting the specific performance claims that can be made.

\section{Why the feedback result matters}
The lineage from the group work to this thesis is explicit. The same spatial sequence and image-construction principles are used, a network is trained at the appropriate duration, and its actual effect on returned candidates is measured. The experiments rule out two trivial explanations for weak utility. The final candidates did not universally collapse to one unusable image, and the selectors did not operate on unequal pools. The score was capable of changing which trajectory a user would receive.

What failed to transfer was the advantage in image-based agreement. Part of a class signal may be tied to location or to how often a trajectory changes grid cells. A participant that looks distinctive to the image model need not be a helpful training example for a network that uses fine timing. This is a hypothesis consistent with the cross-family reversal, rather than a demonstrated decomposition of the image classifier. The fixed-road perturbation in Chapter~\ref{ch:g8} establishes sensitivity to timing, but higher image scores did not select motion closer to the observed class speed distribution.

For the six-class learner, the downstream comparison also changes the distribution of example presentations. Only five classes receive extra traces, every extra trace shares a parent condition, and one generator supplies all augmentation examples. These choices were controlled across selection rules but not between real-only and augmented data. The real-only contrast measures the practical effect of the entire intervention. The random-synthetic contrast more narrowly measures the incremental value of choosing with GAF. Neither comparison supports a general promise of better real-data classification.

\section{Internal and external validity}
The test cohort comes from one simulation export in one geographically aligned scenario. Training, validation and testing are disjoint by participant, but the same network and simulation conventions are shared. Frequent streets can appear in all splits. The experiment evaluates new participants in this setting, rather than transfer to a new map, new demand model or physical vehicle fleet. The exact original SUMO configuration and some participant-to-permission mappings are uncertain.

A second limitation follows from the target representation itself. Monotone progress is well suited to driving forward along an explicit path. It treats lane changes, reversing and complex walking movements poorly. The start-offset correction is an engineering method to preserve the requested first point. The fixed path budget can introduce a cap, even though none occurred in the locked result. Signals, other agents and collision avoidance are outside the conditioning information. These omissions define what a road-consistent single trace means in this experiment.

Statistical uncertainty is dominated by the small number of independent test participants in some classes. Three diffusion fits reduce reliance on one initialization but do not describe the distribution of outcomes across training datasets. The E5 augmentation pool comes from one generator only. Multiple quantities were examined during development and several metrics are reported on the final test. Replica counts, metric-specific intervals and failed controls are therefore retained throughout the report.

\section{Implications for a subsequent study}
The independent simulation in Section~\ref{sec:independent-sumo} supplies a known configuration for new observations. Further replication would benefit from calibrated traffic demand and comparison against the exact original configuration and lane IDs. Those records would separate error introduced by map reconstruction from error introduced by the motion model. Independent runs with different route demand would also permit evaluating how the planner responds to unfamiliar transition frequencies. The current secondary table lacks the matched road provenance and class coverage required for that experiment.

Following the timing intervention of Chapter~\ref{ch:g8}, a complementary study could hold timing fixed while changing a legal route. The change in the image score across these pairs would help distinguish spatial preference from motion sensitivity. A downstream study could also vary the proportion of synthetic examples while holding the number of observed-example presentations constant. Such work would need an independent evaluation set because the present held-out results have been inspected.

\section{Conclusion}
The final product is a full timed trip on a planned route between two endpoints. The clock varies between samples, and a coordinate is returned at the exact sampled arrival. The fixed-window studies isolate behaviour along a local segment. Three controls explain what each component contributes.

The permitted-path decoder accounts for the geometric result. On the 53 eligible locked-test requests, both diffusion and empirical transfer passed the same road audit. Holding those paths fixed, diffusion reduced the discrepancy in travel-distance distributions by 23.7 metres. The evidence for speed and acceleration was less decisive. The restriction to 53 of 270 held-out participants is part of the result.

The diffusion model accounts for a separate, reproducible class response. In development the bus label reduced mean travel by 346 metres relative to the car label. Across 96 new jointly permitted paths the reduction was 343.5 metres, with no change to noise or route within a pair. Secondary interventions showed distinct responses for motorcycle, bicycle and taxi. These comparisons establish control of generated movement on the represented roads. Fidelity to a new traffic regime is a further question, particularly when its bus schedule or demand differs from training.

Scaling the simulator enabled a separate eight-class generator, with independent participant IDs in each training and evaluation episode. At 20,000 updates, adding simulated observations improved travel fit on the G15 population. The G14 primary comparison had a smaller estimate and an interval crossing zero. Seven bicycle samples exhausted their routes, and pedestrian support remained limited to long single sidewalks. The independent terminal-scoring check also left a material benefit from GAF unresolved.

A spatial speed transformation made generated vehicles slow down more realistically on bends, while their steady-moving runs became less faithful. An eight-label implementation of the inherited image feedback frequently changed which candidate was delivered, with no clear gain in either physical pattern on G14 or G15. These trade-offs are visible only when the evaluation preserves temporal order and road context.

The inherited Hilbert--GAF representation helped measure what a recognizer prefers within a candidate pool. A retrained scorer yielded a modest within-request speed benefit relative to the average candidate. That observation did not translate into a reliable improvement in the population-level speed discrepancy or in downstream classification. Fixed-time and state-triggered interventions likewise supplied no stable fidelity benefit on independent traffic. The practical distinction is between requesting a class, recognizing a class and generating its distribution of motion. They require separate measurements, and only the first directly specifies what a synthesis system is asked to do.
